\section{Model and notation}\label{sec-model}

\subsection{Forwards with discrete dividends}
Time is measured in years from the valuation date. Let $S$ be the spot, $r(t)$ the continuously compounded zero
rate used for discounting and funding, and $b(t)$ the zero \emph{repo} (borrow cost) curve, so that holding the
stock financed by a repo earns $r-b$. Write $R(t)=r(t)\,t$, $B(t)=b(t)\,t$ and $G(t)=e^{R(t)-B(t)}$. Cash
dividends $d_i$ go ex at times $t_i$.

\begin{definition}[Forward]\label{def-forward}
The forward for delivery at $T$ is
\begin{equation}\label{eq-forward}
F(T)=G(T)\bigl(S-D(T)\bigr),\qquad D(T)=\sum_{t_i\le T}\frac{d_i}{G(t_i)}.
\end{equation}
$D(T)$ is the \emph{carry-discounted cumulative dividend}: dividends discounted at the carry rate $r-b$.
\end{definition}

\Cref{eq-forward} is the usual replication argument: buy the stock financed at $r-b$, pay the dividends forward.
With flat curves it reads $F(T)=Se^{(r-b)T}-\sum_{t_i\le T}d_ie^{(r-b)(T-t_i)}$. Note that $D$ depends on the repo
curve through $G(t_i)$; this coupling is weak (a few basis points of repo per year on the dividend discounting) but
the bootstrap below handles it exactly.

\subsection{Carry space}
\begin{definition}[Carry]\label{def-carry}
The \emph{cumulative repo} or \emph{carry} is $B(T)=b(T)\,T$. For a quantity observed at pillars $T_1<\dots<T_K$ we
write $B_k=B(T_k)$, $b_k=b(T_k)$, $D_k=D(T_k)$, $F_k=F(T_k)$.
\end{definition}

\begin{lemma}[Additivity in carry]\label{lem-additivity}
$\ln F(T)=R(T)-B(T)+\ln\bigl(S-D(T)\bigr)$. Consequently, perturbations of the forward, the dividends and the carry
satisfy, to first order,
\begin{equation}\label{eq-additivity}
\delta B(T)=-\,\delta\ln F(T)-\frac{\delta D(T)}{S-D(T)}+\delta R(T).
\end{equation}
\end{lemma}
\begin{proof}
Take logarithms in \cref{eq-forward}; differentiate, using $\delta\ln(S-D)=-\delta D/(S-D)+O(\delta D^2)$.
\end{proof}

\begin{insight}
In carry space the mechanisms add up and nothing is amplified. A carry error $\varepsilon$ at tenor $T$ is a repo
error $\varepsilon/T$: the same $0.1\%$ error in the forward is $120$\,bp of repo at one month and $2$\,bp at five
years. This is why implied repo curves look so noisy at the short end, and why all tests in this note work on $B$,
not on $b$.
\end{insight}

\subsection{The implied repo map}
Given forwards $F_1,\dots,F_K$ at the pillars, the rate curve and the dividends, the \emph{implied repo} is the repo
curve, piecewise linear in $B$ (flat forward repo between pillars, $B(0)=0$), that reprices the forwards. It is
solved pillar by pillar. Fix $B_1,\dots,B_{k-1}$ and let $x=B_k$. For a dividend between the previous pillar and
$T_k$, $B(t_i)=B_{k-1}+w_i(x-B_{k-1})$ with $w_i=(t_i-T_{k-1})/(T_k-T_{k-1})\in(0,1]$, so the model forward at $T_k$ is
\begin{equation}\label{eq-phik}
\Phi_k(x)=Se^{R_k-x}-\sum_{t_i\le T_{k-1}}d_i\,e^{(R_k-R_i)-(x-B(t_i))}
-\sum_{T_{k-1}<t_i\le T_k}d_i\,e^{(R_k-R_i)-(1-w_i)(x-B_{k-1})}.
\end{equation}

\begin{lemma}[Monotonicity]\label{lem-monotone}
$\Phi_k'(x)\le-\Phi_k(x)$ for all $x$. In particular $\Phi_k$ is strictly decreasing wherever the forward is positive,
and $g_k(x)=\ln\Phi_k(x)-\ln F_k$ satisfies $g_k'(x)\le-1$.
\end{lemma}
\begin{proof}
Differentiating \cref{eq-phik} term by term,
$\Phi_k'(x)=-Se^{R_k-x}+\sum_{t_i\le T_{k-1}}d_ie^{(\cdot)}+\sum_{T_{k-1}<t_i\le T_k}(1-w_i)\,d_ie^{(\cdot)}$.
Every dividend term is positive and $1-w_i\le1$, so replacing $(1-w_i)$ by $1$ can only increase the right-hand
side, which then equals $-\Phi_k(x)$. Dividing by $\Phi_k>0$ gives $g_k'\le-1$.
\end{proof}

\begin{proposition}[Well-posed bootstrap]\label{prop-bootstrap}
Assume $S-D(T_{k-1})>0$ for the dividends already solved. Then for every $F_k>0$ there is a unique $B_k$ with
$\Phi_k(B_k)=F_k$ and $\Phi_k>0$ on the solution branch. Newton's method on $g_k$ converges locally quadratically,
and globally when safeguarded by bisection.
\end{proposition}
\begin{proof}
As $x\to-\infty$,
\[
e^{x}\Phi_k(x)\to e^{R_k}\Bigl(S-\sum_{t_i\le T_{k-1}}d_ie^{-R_i+B(t_i)}\Bigr)=e^{R_k}\bigl(S-D(T_{k-1})\bigr)>0,
\]
because the in-period terms grow only like $e^{-(1-w_i)x}$; hence $\Phi_k\to+\infty$. As $x$ increases, $\Phi_k$
decreases strictly while positive (\cref{lem-monotone}) and either tends to $0$ or crosses zero, so by continuity
it takes every value in $(0,\infty)$ exactly once on the branch where it is positive. On that branch $g_k$ is
$C^2$ with $|g_k'|\ge1$, so the Newton iteration $x\mapsto x-g_k(x)/g_k'(x)$ converges quadratically near the root;
the bisection safeguard (keep a bracket $[\ell,u]$ with $g_k(\ell)>0>g_k(u)$ and fall back to its midpoint) makes
convergence global.
\end{proof}

\begin{example}[A bootstrap by hand]\label{ex-bootstrap}
Spot $100$, flat rate $3\%$, dividends $1.5$ at $t=0.30$ and $2.0$ at $t=0.80$, pillars $0.25$, $0.5$, $1$, true
repo $80$, $70$, $60$\,bp. The forwards and the implied repo are
\begin{center}\input{generated/\lang/tab_boot}\end{center}
At the third pillar the Newton residual $|g_3|$ falls from \vNewtonA{} to \vNewtonB{} in one step and to zero in the
next. Without dividends $g_k$ is linear and a single step suffices.
\end{example}

\subsection{Forwards from consensus option prices}
\begin{lemma}[Parity and the consensus forward]\label{lem-parity}
Let each contributor $j$ quote European calls and puts with a common discount factor $DF$. Then
\begin{enumerate}[label=(\roman*)]
\item $C_j(K)-P_j(K)=DF\,(F_j-K)$ for every strike;
\item if the consensus averages the same contributors at every strike, $\overline{C}(K)-\overline{P}(K)=DF\,(\overline{F}-K)$
with $\overline{F}$ the mean of the contributors' forwards;
\item the strike $K^\ast$ where $\overline{C}=\overline{P}$ equals $\overline{F}$ and does not depend on $DF$;
multiplying all prices by a common factor leaves $K^\ast$ unchanged.
\end{enumerate}
\end{lemma}
\begin{proof}
(i) is put--call parity~\cite{hull}. (ii) follows by averaging (i), which is linear. (iii) Set (ii) to zero; scaling both
prices scales the line $K\mapsto\overline{C}-\overline{P}$ without moving its root.
\end{proof}

\begin{example}\label{ex-parity}
With strikes $90$, $100$, $110$, puts $2.1$, $5.4$, $10.9$ and calls built from $F=102.3$, $DF=0.97$, the weighted
regression of $\overline{C}-\overline{P}$ on $K$ returns $F=\vParF$ and $DF=\vParDF$. Scaling all prices by $1.02$
(a wrong discount factor) returns $F=\vParFtwo$ and $DF=\vParDFtwo$: the forward is unchanged.
\end{example}

\begin{remark}
Two cases break (ii)--(iii): outlier screening applied separately per strike and per option type (the averaged
contributors differ between calls and puts), and American exercise (parity becomes an inequality). The fitted $DF$
compared with the curve gives an implied funding spread that flags such expiries. The controller keeps its own
rate curve on both sides of the comparison, so rates cancel in the Totem--FO difference.
\end{remark}

\subsection{Dividend buckets and allocation}
\begin{definition}[Allocation]\label{def-allocation}
Totem dividend levels $L_y$ are sums of dividends going ex in period $y=(e_{y-1},e_y]$ (typically calendar years). A
profile assigns weights $\pi_i\ge0$ with $\sum_{i\in y}\pi_i=1$ to the ex-dates of each period; the allocated
dividends are $d_i=L_y\pi_i$. The controller normally takes the FO ex-date calendar as profile.
\end{definition}

\begin{lemma}\label{lem-allocation}
Holding $G$ fixed, $D_k=\sum_yA_{ky}L_y$ with $A_{ky}=\sum_{i\in y,\,t_i\le T_k}\pi_i/G(t_i)$.
\end{lemma}
\begin{proof}
Substitute $d_i=L_y\pi_i$ in \cref{eq-forward} and group by period.
\end{proof}
$A_{ky}$ is the carry-discounted share of period $y$ that goes ex before expiry $k$. It is $0$ or close to $1/G$
except for the periods that straddle an expiry, which is exactly where the profile matters.

\subsection{Totem data model}
\begin{assumption}[Consistent contributors]\label{ass-consistent}
Contributor $j$ holds views $b_j(\cdot)$ and $D_j(\cdot)$ and quotes option prices consistent with
$F_j(T)=G_j(T)(S-D_j(T))$, with the common rate curve and spot.
\end{assumption}
\begin{assumption}[Panels]\label{ass-panels}
At expiry $T_k$ the option panel $\Aset_k$ has $n_{A,k}$ contributors; for period $y$ the dividend panel $\Bset_y$
has $n_{B,y}$. Panels are nested in tenor ($\Aset_{k+1}\subseteq\Aset_k$, $\Bset_{y+1}\subseteq\Bset_y$). View errors
are independent across contributors. The two services share $m$ contributors ($m=0$ for disjoint panels).
\end{assumption}
\begin{assumption}[Consensus]\label{ass-consensus}
The consensus is the arithmetic mean over the panel. Totem reports, per expiry and per period, the consensus, the
dispersion across contributors and their number.
\end{assumption}

\begin{definition}[Totem-implied repo and carry difference]\label{def-implied}
The Totem-implied repo $b^\T$ is the bootstrap of the consensus forwards $\overline{F}_k$ with the allocated
consensus dividends $\overline{L}_y$. With the FO repo $b^\FO$ at the same pillars, the \emph{carry difference} is
\begin{equation}\label{eq-y}
y_k=B^\T_k-B^\FO_k=\bigl(b^\T_k-b^\FO_k\bigr)\,T_k,\qquad k=1,\dots,K.
\end{equation}
\end{definition}
